\section{The \texorpdfstring{$\Phi_3$}{Phi3} Exclusion Criterion}\label{sec:phi3}

In this section we develop a structural obstruction to degree-one capture
in the FF-EM sequence that has no analog over $\ZZ$. The result is
unconditional and gives a concrete family of primes $p$ for which the
autonomous map $f(w) = w(w+1)$ \emph{cannot} reach the death residue $-1$.

\subsection{The autonomous map and degree-one targets}

Recall from Definition~\ref{def:autonomous} that the autonomous map
$f \colon \Fp \to \Fp$ is given by $f(w) = w(w+1)$. For a degree-one
target $Q = (t - a)$, the walk $\ffprod(n) \bmod Q$ in $\Fp$ follows
the recursion $w_{n+1} = f(w_n)$ whenever the walk multiplier
$\ffseq(n+1) \bmod Q$ equals $w_n + 1$ (the self-consistent case).

The target $Q$ divides $\ffprod(n) + 1$ if and only if
$\ffprod(n) \equiv -1 \pmod{Q}$, i.e., the walk position $w_n = -1$
in $\Fp$. Hence the question of degree-one capture reduces to:
\emph{does the orbit of $f$ reach $-1$?}

\subsection{Fixed points and absorption}

\begin{lemma}[Absorption at zero]\label{lem:absorption}
\leancheck{ffAutonomousMap\_neg\_one}
\leancheck{ffAutonomousMap\_zero}
The map $f$ satisfies $f(-1) = 0$ and $f(0) = 0$. In particular, $0$ is
an absorbing fixed point: once the orbit reaches $0$, it remains there
forever.
\end{lemma}

\begin{proof}
$f(-1) = (-1)(0) = 0$ and $f(0) = 0 \cdot 1 = 0$. If $w_n = 0$ for
some $n$, then $w_m = 0$ for all $m \geq n$ by induction.
\end{proof}

\leancheck{ffAutonomousOrbit\_absorbed}

\subsection{The preimage of \texorpdfstring{$-1$}{-1}}

\begin{lemma}[$\Phi_3$ criterion]\label{lem:phi3-criterion}
\leancheck{ffAutonomousMap\_eq\_neg\_one\_iff}
For $w \in \Fp$,
\[
  f(w) = -1 \quad \Longleftrightarrow \quad w^2 + w + 1 = 0.
\]
That is, the preimage of $-1$ under $f$ is exactly the root set of the
third cyclotomic polynomial $\Phi_3(w) = w^2 + w + 1$.
\end{lemma}

\begin{proof}
$f(w) = w(w+1) = w^2 + w$. Hence $f(w) = -1$ iff $w^2 + w = -1$ iff
$w^2 + w + 1 = 0$.
\end{proof}

\subsection{Absence of cube roots of unity}

The roots of $\Phi_3$ are the primitive cube roots of unity. Whether
$\Phi_3$ has roots in $\Fp$ depends on $p \bmod 3$.

\begin{lemma}\label{lem:cube-root}
\leancheck{pow\_three\_eq\_one\_of\_phi3}
If $w^2 + w + 1 = 0$ in $\Fp$, then $w^3 = 1$.
\end{lemma}

\begin{proof}
$(w-1)(w^2+w+1) = w^3 - 1$, so $w^2 + w + 1 = 0$ gives $w^3 = 1$.
\end{proof}

\begin{theorem}[No roots of $\Phi_3$ for $p \equiv 2 \pmod{3}$]\label{thm:phi3-no-roots}
\leancheck{phi3\_no\_roots}
If $p \equiv 2 \pmod{3}$ and $p \geq 5$, then $w^2 + w + 1 \neq 0$ for
all $w \in \Fp$.
\end{theorem}

\begin{proof}
Suppose $w^2 + w + 1 = 0$ for some $w \in \Fp$.

\emph{Step 1}: $w^3 = 1$ (Lemma~\ref{lem:cube-root}).

\emph{Step 2}: $w \neq 0$ (substituting $w = 0$ gives $1 = 0$, a
contradiction since $p \geq 2$).

\emph{Step 3}: $w \neq 1$ (substituting $w = 1$ gives $3 = 0$ in $\Fp$;
since $p \geq 5$, we have $p \nmid 3$, a contradiction).

\emph{Step 4}: $w$ is a unit in $\Fp^\times$ (since $w \neq 0$ in a field).
Since $w^3 = 1$ and $w \neq 1$, the order of $w$ in $\Fp^\times$ is
exactly $3$.

\emph{Step 5}: By Lagrange's theorem, $3 \mid |\Fp^\times| = p - 1$.

\emph{Step 6}: But $p \equiv 2 \pmod{3}$ gives $p - 1 \equiv 1 \pmod{3}$,
so $3 \nmid (p-1)$. Contradiction.
\end{proof}

\begin{remark}
The condition $p \geq 5$ excludes $p = 2$ (where $\Phi_3$ is irreducible:
$1 + 1 + 1 = 1 \neq 0$ in $\mathbb{F}_2$, and $0 + 0 + 1 = 1 \neq 0$)
and $p = 3$ (where $w^2 + w + 1 = (w-1)^2$ in $\mathbb{F}_3$, so $w = 1$
is a double root---but $p = 3$ has $3 \bmod 3 = 0$, not $2$).
The primes satisfying $p \equiv 2 \pmod{3}$ include
$5, 11, 17, 23, 29, 41, 47, \ldots$
\end{remark}

\subsection{The main exclusion theorem}

\begin{theorem}[Unreachability of $-1$]\label{thm:unreachable}
\leancheck{ff\_neg\_one\_unreachable}
For $p \equiv 2 \pmod{3}$ with $p \geq 5$: the orbit of $f$ from any
starting value $a \neq -1$ never hits $-1$. That is,
\[
  a \neq -1 \implies f^n(a) \neq -1 \quad \text{for all } n \geq 0.
\]
\end{theorem}

\begin{proof}
By induction on $n$. The base case $n = 0$ is the hypothesis $a \neq -1$.
For $n + 1$: $f^{n+1}(a) = f(f^n(a))$, and by
Theorem~\ref{thm:phi3-no-roots}, $f(w) \neq -1$ for \emph{all} $w \in \Fp$
(since $\Phi_3$ has no roots). Hence $f^{n+1}(a) \neq -1$.
\end{proof}

\begin{corollary}\label{cor:visited-once}
\leancheck{ff\_neg\_one\_visited\_at\_most\_once}
For $p \equiv 2 \pmod{3}$ with $p \geq 5$: the value $-1$ is visited by
the orbit at most at step $0$ (i.e., only if $a = -1$). At every step
$n \geq 1$, the orbit value is never $-1$, regardless of the starting value.
\end{corollary}

\begin{proof}
For $n \geq 1$, the orbit value $f^n(a)$ is in the range of $f$. But
$-1 \notin \mathrm{range}(f)$ by Theorem~\ref{thm:phi3-no-roots}.
\end{proof}

\subsection{Interpretation: a structural obstruction}

Theorem~\ref{thm:unreachable} provides a \emph{structural obstruction} to
degree-one capture in the FF-EM sequence under the autonomous map model.
For primes $p \equiv 2 \pmod{3}$ with $p \geq 5$:

\begin{itemize}
  \item If the walk position $\ffprod(0) \bmod Q$ is $-1$, then $Q$ divides
    $\ffprod(0) + 1 = t + 1$, and $Q$ may be captured at step~$1$.
  \item If $\ffprod(0) \bmod Q \neq -1$, the walk \emph{never} reaches $-1$
    under the autonomous map, so $Q$ is never captured via the degree-one
    route.
\end{itemize}

This gives a concrete family of targets that are unreachable by the
autonomous map, constituting a degree-one instance of the orbit-specificity
barrier.

\subsection{Comparison with the integer case}

Over $\ZZ$, only four primes are known to be excluded by elementary methods
from certain positions in the Euclid--Mullin sequence. The function field
$\Phi_3$ criterion is \emph{stronger}: it excludes $p - 2$ out of $p - 1$
possible starting residues from ever reaching $-1$ (only $a = -1$ itself
can reach $-1$).

However, the $\Phi_3$ criterion applies only under the ``perpetual
irreducibility'' assumption (that the walk multiplier is determined by
the autonomous map). In practice, the walk multiplier $\ffseq(n+1)
\bmod Q$ need not equal $w_n + 1$; it is determined by the specific
factorization of $\ffprod(n) + 1$, which can break the autonomous map
structure.

\begin{remark}
The $\Phi_3$ criterion is specific to the function field setting because
the autonomous map $f(w) = w(w+1)$ arises from the polynomial structure
of $\Fpt$. Over $\ZZ$, the ``walk'' modulo $q$ is $P(n) \bmod q$, but
the multiplier $a_{n+1} \bmod q$ is not determined by $P(n) \bmod q$
alone (it depends on the global factorization of $P(n)+1$, not just on
the residue class).
\end{remark}

\subsection{The cyclotomic structure: Dead End \#129}

The $\Phi_3$ criterion connects to a deeper structural phenomenon. Over
$\mathbb{F}_2$, the product of all monic irreducibles of degree dividing $d$
equals $t^{2^d} - t$, so adding $1$ gives the polynomial $t^{2^d} - t + 1$.
For $d = 2$:
\[
  t(t+1)(t^2+t+1) = t^4 + t, \qquad
  t^4 + t + 1 = \Phi_5(t),
\]
the $5$th cyclotomic polynomial. Its Galois group is
$\mathrm{Gal}(\mathbb{F}_{2^4}/\mathbb{F}_2) \cong \ZZ/4\ZZ$, which is
abelian, \emph{not} $S_4$.

More generally, the classical identity
$\prod_{\deg Q \mid d} Q = t^{p^d} - t$ forces the greedy FF-EM sequence
to build products dividing $t^{p^d} - t$, so the ``$+1$ shifted'' polynomial
$\ffprod(n) + 1$ has its splitting field contained in $\GF{p}{N}$ for some
$N$. The Galois group is therefore \emph{cyclic} (hence abelian), generated
by the Frobenius automorphism.

This means the ``large monodromy'' hypothesis ($\mathrm{FFLM}$: the Galois
group eventually contains the alternating group $A_d$) is
\emph{structurally false} for the greedy FF-EM sequence. The abelian Galois
group recovers only the Weil bound information (Dirichlet characters
correspond to abelian representations), which is population equidistribution
---the same information we already have.

\leancheck{ff\_cyclotomic\_dead\_end}
